% ===== PRÉAMBULE CANONIQUE — à coller VERBATIM en tête de CHAQUE .tex =====
\documentclass[11pt,a4paper]{article}
\usepackage[utf8]{inputenc}\usepackage[T1]{fontenc}\usepackage{lmodern}
\usepackage[french]{babel}
\usepackage{amsmath,amssymb,mathtools}\usepackage{icomma}
\usepackage{siunitx}\sisetup{locale=FR}  % \qty{}{}, \si{}, \ang{}
\usepackage{xcolor}\usepackage{colortbl}
\usepackage{tikz}\usepackage{pgfplots}\pgfplotsset{compat=1.18}
\usepackage[siunitx,europeanresistors]{circuitikz} % circuits (élec.)
\usepackage[most]{tcolorbox}\usepackage{enumitem}\usepackage{array,booktabs}
\usepackage[a4paper,margin=2.2cm]{geometry}
\usetikzlibrary{arrows.meta,calc,angles,quotes,positioning,patterns,%
  backgrounds,shapes.geometric,decorations.pathmorphing,decorations.markings,intersections}
\definecolor{primary}{RGB}{222,49,99}\definecolor{primarydark}{RGB}{150,22,55}
\definecolor{defcol}{RGB}{0,114,178}\definecolor{methcol}{RGB}{52,92,148}
\definecolor{excol}{RGB}{0,135,90}\definecolor{attncol}{RGB}{200,40,55}
\tcbset{boxbase/.style={enhanced,breakable,boxrule=0.9pt,arc=2.5pt,
  left=3mm,right=3mm,top=2.3mm,bottom=2.3mm,fonttitle=\bfseries,coltitle=white}}
% --- boîtes TUTORIEL (noms EXACTS, ne pas renommer) ---
\newtcolorbox{cle}[1][]{boxbase,colback=defcol!6,colframe=defcol,
  title={\textsc{À retenir}\ifx\relax#1\relax\relax\else\textmd{ : \itshape #1}\fi}}
\newtcolorbox{methode}[1][]{boxbase,colback=methcol!7,colframe=methcol,sharp corners,
  title={\textsc{Méthode}\ifx\relax#1\relax\relax\else\textmd{ --- \itshape #1}\fi}}
\newtcolorbox{exemplebox}[1][]{boxbase,colback=excol!5,colframe=excol,
  title={\textsc{Exemple}\ifx\relax#1\relax\relax\else\textmd{ : \itshape #1}\fi}}
\newtcolorbox{piege}[1][]{boxbase,colback=attncol!6,colframe=attncol,
  title={\textsc{Piège}\ifx\relax#1\relax\relax\else\textmd{ : \itshape #1}\fi}}
\newtcolorbox{remarque}[1][]{boxbase,colback=black!4,colframe=black!45,
  title={\textsc{Remarque}\ifx\relax#1\relax\relax\else\textmd{ : \itshape #1}\fi}}
% --- boîtes EXERCICE / CORRIGÉ (noms EXACTS, auto-comptées) ---
\newtcolorbox[auto counter]{exo}[1][]{boxbase,colback=primary!4,colframe=primary,
  title={\textsc{Exercice}~\thetcbcounter\ifx\relax#1\relax\relax\else\textmd{ --- \itshape #1}\fi}}
\newtcolorbox[auto counter]{corrige}[1][]{boxbase,colback=excol!4,colframe=excol,
  title={\textsc{Corrigé}~\thetcbcounter\ifx\relax#1\relax\relax\else\textmd{ --- \itshape #1}\fi}}
\newcommand{\vv}[1]{\vec{#1}}\newcommand{\ds}{\displaystyle}
\newcommand{\R}{\mathbb{R}}\newcommand{\N}{\mathbb{N}}\newcommand{\Z}{\mathbb{Z}}
% =========================================================================
\begin{document}
% THEME_TITLE: La résistance d'un conducteur : loi de Pouillet et température
\usetikzlibrary{babel}

\begin{center}
{\LARGE\bfseries\color{primarydark} Thème 3 — Résistance d'un fil : loi de Pouillet et température}\\[2pt]
{\color{primary}\itshape $R=\rho\,L/S$, résistivité, conductivité et effet thermique}
\end{center}

\begin{cle}[loi de Pouillet]
La résistance d'un fil conducteur homogène de section constante vaut :
\[ \boxed{\;R=\rho\,\dfrac{L}{S}\;} \]
\begin{itemize}[leftmargin=1.5em,topsep=1pt,itemsep=0pt]
  \item $R$ en ohms (\si{\ohm}) ; \quad $L$ = longueur en \si{\metre} ; \quad $S$ = aire de la section
        en \si{\metre\squared} ;
  \item $\rho$ = \textbf{résistivité} du matériau, en \si{\ohm\metre} (caractérise sa \emph{nature}).
\end{itemize}
\end{cle}

\begin{center}
\begin{tikzpicture}[scale=1,>=Stealth]
  \draw[fill=primary!15,draw=primarydark] (0,0) rectangle (5.2,0.7);
  \draw[<->] (0,-0.28) -- (5.2,-0.28) node[midway,below,font=\footnotesize]{longueur $L$};
  \draw[->] (5.55,0.35) -- (6.1,0.35);
  \draw[primarydark,fill=primary!15] (6.35,0.35) ellipse (0.18 and 0.35);
  \node[right,font=\footnotesize] at (6.6,0.35){section $S$};
  \draw[<->] (-0.28,0) -- (-0.28,0.7) node[midway,left,font=\scriptsize]{};
\end{tikzpicture}
\end{center}

\begin{cle}[proportionnalités à retenir]
$R$ est \textbf{proportionnelle à $L$} et \textbf{inversement proportionnelle à $S$} :
\begin{itemize}[leftmargin=1.5em,topsep=1pt,itemsep=0pt]
  \item on \emph{double} la longueur $\Rightarrow$ $R$ \emph{double} ;
  \item on \emph{double} la section $\Rightarrow$ $R$ est \emph{divisée par 2} ;
  \item un \emph{bon} conducteur a une $\rho$ \emph{petite} (cuivre, argent) ; un \emph{isolant} a une
        $\rho$ \emph{énorme} (verre).
\end{itemize}
\end{cle}

\begin{cle}[conductivité et quelques résistivités]
La \textbf{conductivité} est l'inverse de la résistivité : $\sigma=\dfrac{1}{\rho}$ (en
\si{\siemens\per\metre}). Plus $\sigma$ est grande, meilleur est le conducteur.
\begin{center}\footnotesize
\renewcommand{\arraystretch}{1.1}
\begin{tabular}{lcc}
\toprule
\textbf{Matériau} & $\rho$ (\si{\ohm\metre}) & \textbf{type} \\
\midrule
Argent & $1{,}6\times10^{-8}$ & excellent conducteur \\
Cuivre & $1{,}7\times10^{-8}$ & excellent conducteur \\
Aluminium & $2{,}8\times10^{-8}$ & bon conducteur \\
Fer & $10\times10^{-8}$ & conducteur \\
Verre & $10^{12}$ à $10^{16}$ & isolant \\
\bottomrule
\end{tabular}
\end{center}
\end{cle}

\begin{cle}[effet de la température (métal pur)]
Pour un métal pur, la résistivité (donc la résistance) croît \emph{linéairement} avec la température :
\[ \boxed{\;\rho=\rho_0\,(1+\alpha T)\;}\qquad R=R_0\,(1+\alpha T) \]
$\rho_0$, $R_0$ = valeurs à $T=0\,\si{\degreeCelsius}$ ; $\alpha$ = coefficient de température
(\si{\per\degreeCelsius}). Valable si $T$ n'est pas trop élevée.
\end{cle}

\begin{methode}[appliquer la loi de Pouillet]
\begin{enumerate}[leftmargin=1.7em,topsep=1pt,itemsep=0pt]
  \item \textbf{Convertir la section} en \si{\metre\squared} : $\qty{1}{\milli\metre\squared}=10^{-6}\,\si{\metre\squared}$.
  \item Appliquer $R=\rho L/S$. Pour un \emph{rapport} (doubler $L$, etc.), inutile de tout recalculer.
  \item Température : passer par le rapport $\dfrac{R(T)}{R_0}=1+\alpha T$.
\end{enumerate}
\end{methode}

\begin{piege}[deux pièges d'unités]
La section est souvent donnée en $\si{\milli\metre\squared}$ : ne pas oublier de la convertir en
$\si{\metre\squared}$ ($\times10^{-6}$). Et $\alpha T$ est \emph{sans dimension} : $\alpha$ en
$\si{\per\degreeCelsius}$ multiplie $T$ en $\si{\degreeCelsius}$.
\end{piege}

\end{document}
